\documentclass[11pt]{article}
\usepackage{../../latex/stat4ml_style}

\title{Chapter 10: Uncertainty and Modern Generative Models \\ \large Part III --- Statistical Foundations of Deep Learning \& Reinforcement Learning}
\author{Omid Safarzadeh}
\date{}

\begin{document}
\maketitle

\begin{abstract}
\noindent
A model that only predicts is half a model; the other half is knowing when not to
trust it, and generating new data from what it learned. This closing chapter
gathers three statistical ideas that give modern deep learning both. Conformal
prediction turns any model into one with a \emph{finite-sample, distribution-free}
coverage guarantee --- valid prediction sets with no assumption about the data or
the model. Calibration asks the sharper question of whether a classifier's
probabilities mean what they say, and measures the gap. And score matching ---
learning the gradient of the log-density --- is the statistical engine inside
diffusion models, which generate by learning to reverse the addition of noise.
Each is a clean piece of statistics doing work people now take for granted.
\end{abstract}

\tableofcontents

% ======================================================================
\section{Conformal prediction}
% ======================================================================

A point prediction says nothing about its own reliability. \keyterm{Conformal
prediction} wraps any predictor in a set that is guaranteed to contain the truth
a chosen fraction of the time --- with no distributional assumptions and for any
sample size.

\begin{definition}[Split conformal prediction]
Split the data into a training part and a \keyterm{calibration} part. Fit any
predictor $\hat f$ on the training part. On the calibration part compute
\keyterm{nonconformity scores} $s_i = |y_i - \hat f(x_i)|$, and let $q$ be the
$\lceil (n+1)(1-\alpha)\rceil / n$ empirical quantile of the $s_i$. The prediction
set for a new $x$ is $[\hat f(x)-q,\ \hat f(x)+q]$.
\end{definition}

\begin{theorem}[Marginal coverage]
If the calibration and test points are exchangeable, the conformal set covers the
true $y$ with probability at least $1-\alpha$:
$\Pr\big(y\in[\hat f(x)-q,\hat f(x)+q]\big) \ge 1-\alpha$.
\end{theorem}
\begin{proof}[Proof idea]
By exchangeability the test score $s_{\text{new}}=|y-\hat f(x)|$ is equally likely
to land in any position among the $n+1$ scores $\{s_1,\dots,s_n,s_{\text{new}}\}$.
The set fails only when $s_{\text{new}}$ exceeds the chosen quantile, which happens
for at most $\lceil (n+1)(1-\alpha)\rceil$ of the $n+1$ equally likely ranks,
i.e.\ with probability at most $\alpha$.
\end{proof}

The guarantee leans on nothing but exchangeability --- not on $\hat f$ being
correct, nor on the noise being Gaussian. The left panel shows a $90\%$ band built
this way; measured on fresh data it covers $90.4\%$, exactly the promised rate.
The band's \emph{width} reflects the model's accuracy (a good $\hat f$ gives tight
bands), but its \emph{coverage} is guaranteed regardless.

% ======================================================================
\section{Calibration}
% ======================================================================

A classifier that outputs $0.9$ should be right $90\%$ of the time on the inputs
it is that confident about. When this holds the model is \keyterm{calibrated}; when
it does not, its probabilities are misleading even if its top-1 accuracy is good
(modern networks are often badly overconfident).

\begin{definition}[Calibration and ECE]
A model is \keyterm{calibrated} if $\Pr(y=1\mid \hat p = p)=p$ for all $p$. The
\keyterm{expected calibration error} bins predictions by confidence and averages
the gap between confidence and accuracy,
\[
  \mathrm{ECE} = \sum_{b} \frac{|B_b|}{N}\,\big|\operatorname{acc}(B_b) - \operatorname{conf}(B_b)\big|.
\]
\end{definition}

The \keyterm{reliability diagram} plots accuracy against confidence per bin; a
calibrated model lies on the diagonal, and ECE is its average vertical distance
from it. The middle panel shows a calibrated model tracking the diagonal and an
overconfident one bowing away from it. Calibration is a proper-scoring-rule story
(Chapter~5): cross-entropy is minimised by the true probabilities, so a
well-trained model should be calibrated --- that it often is not, after modern
training, is why post-hoc recalibration (temperature scaling) is standard.

% ======================================================================
\section{Score matching and diffusion}
% ======================================================================

To \emph{generate} data you need its distribution $p(x)$, but the normalising
constant of a flexible $p$ is intractable. The escape is to learn not $p$ but its
\keyterm{score}, the gradient of the log-density $\nabla_x\log p(x)$ --- which does
not involve the normaliser at all, since the constant vanishes under the
gradient.

\begin{definition}[Denoising score matching]
Add Gaussian noise, $\tilde x = x + \sigma\varepsilon$ with
$\varepsilon\sim N(0,I)$, and train a network $s_\theta$ to predict the score of
the noised data by regressing on the noise:
\[
  \min_\theta\ \mathbb{E}\big\|\,s_\theta(\tilde x) - \nabla_{\tilde x}\log p_\sigma(\tilde x\mid x)\,\big\|^2,
  \qquad \nabla_{\tilde x}\log p_\sigma(\tilde x\mid x) = -\frac{\tilde x - x}{\sigma^2}.
\]
Learning the score reduces to the supervised problem of denoising.
\end{definition}

The score is a vector field pointing toward higher density, as the right panel
shows for a two-mode distribution: the arrows flow into the modes. A
\keyterm{diffusion model} uses this twice over. A fixed forward process gradually
turns data into pure noise; the model learns the score at every noise level; and
to \keyterm{generate}, it starts from noise and follows the learned score
``uphill'', step by step, reversing the diffusion back to a sample from $p$.
Generation is gradient ascent on a learned log-density --- the same score function
that ran through maximum likelihood (Chapter~2) and the policy gradient
(Chapter~7), now driving a sampler. The course ends where it is fitting: the
score, one of statistics' oldest objects, at the centre of its newest models.

\begin{summarybox}
Conformal prediction wraps any model in prediction sets with a finite-sample,
distribution-free coverage guarantee $\ge 1-\alpha$, resting only on
exchangeability: calibrate nonconformity scores and use their
$(1-\alpha)$-quantile. Calibration asks whether predicted probabilities match
empirical frequencies; the reliability diagram and ECE measure the gap, and
proper scoring rules (Chapter~5) say a well-fit model should be calibrated.
Generative modelling learns the score $\nabla_x\log p$ (no normaliser needed) by
denoising score matching; a diffusion model learns the score across noise levels
and generates by following it from noise back to data --- gradient ascent on a
learned log-density, the score function of Chapters~2 and~7 at the heart of modern
generative AI.
\end{summarybox}

\end{document}
